\subsection{Projektplan}

\todo{Better title; Introductory sentence}

\subsubsection{Marktsituation}

Singleplayer-Kartenspiele und rundenbasierte Roguelites haben sich in den vergangenen Jahren als etabliertes und wirtschaftlich relevantes Segment des Indie-Computerspielmarktes entwickelt.
Insbesondere Spiele, die strategische Entscheidungen mit einer Run-basierten Struktur, wechselnden Spielstrategien und hoher Wiederspielbarkeit verbinden, haben eine große und langfristig aktive Zielgruppe hervorgebracht.

Das hohe Marktpotenzial dieses Segments zeigen zunächst einige außergewöhnlich erfolgreiche Vertreter.
Einer der frühen und genreprägenden Titel ist \emph{Slay the Spire}, das bereits 2019 mehr als 1,5 Millionen verkaufte Einheiten erreicht hatte.
In den folgenden Jahren haben weitere Spiele gezeigt, dass dabei nicht nur direkte Nachfolger dieses Spielprinzips erfolgreich sein können.
\emph{Inscryption} verband Kartenspielmechaniken mit Puzzle-, Horror- und narrativen Elementen und wurde innerhalb von weniger als drei Monaten mehr als eine Million Mal gekauft.
\emph{Balatro} interpretierte das strategische Singleplayer-Kartenspiel erneut auf grundlegend andere Weise und erreichte innerhalb seines ersten Jahres mehr als fünf Millionen Verkäufe.
Dass das Interesse an diesem Marktsegment weiterhin sehr hoch ist, zeigt zuletzt \emph{Slay the Spire 2}: Bereits in der ersten Woche seines Early-Access-Starts im März 2026 wurden drei Millionen Exemplare verkauft und mehr als 25 Millionen Runs gespielt.

Diese außergewöhnlichen Millionenerfolge bilden jedoch nur die besonders sichtbare Spitze eines deutlich breiteren Marktes.
Neben ihnen existiert eine große Zahl kleinerer und mittlerer Produktionen im Bereich der Singleplayer-Kartenspiele und rundenbasierten Roguelites, die innerhalb des Genres eine stabile Zielgruppe erreichen und wirtschaftlich erfolgreich sind.
Der Markt wird damit nicht allein von einzelnen Ausnahmeerscheinungen getragen, sondern von einer etablierten Spielerschaft mit kontinuierlichem Intereasse an neuen Titeln und neuen Interpretationen dieser Spielprinzipien.

Gerade die Vielfalt erfolgreicher Spiele innerhalb dieses Umfelds ist für \coc\ relevant.
Titel wie \emph{Slay the Spire} stehen für klassisches Deckbuilding, während \emph{Balatro} oder \emph{Inscryption} zeigen, dass sich dieselbe grundsätzliche Zielgruppe auch mit deutlich anderen Karten- und Spielsystemen erreichen lässt.
Gleichzeitig demonstrieren Roguelites wie \emph{Hades} oder \emph{Shogun Showdown} die Attraktivität eines anderen Ansatzes zur Entwicklung der Spielstrategie: Spieler ziehen ihre Handlungsmöglichkeiten nicht fortlaufend zufällig aus einem Deck, sondern stellen während eines Runs gezielt eine Auswahl an Fähigkeiten, Verbesserungen und Synergien zusammen.
% Das ist jetzt nicht so richtig ein *Gegensatz*, dass passiert in Deckbuildern halt auf der höheren Ebene.
% Insbesondere ist zu sagen (aber nicht hier in diesem Antrag), dass ein Roguelite ohne RNG in den Upgrades halt auch total lame ist, weil man dann doch einfach das gleiche wieder spielt.

\coc\ liegt an der Schnittstelle dieser Zielgruppen.
Das Spiel richtet sich insbesondere an Spieler strategischer Singleplayer-Kartenspiele und rundenbasierter Roguelites, spricht durch sein Loadout-Building aber ebenso Spieler an, die aus anderen Roguelites die gezielte Entwicklung individueller Archetypen oder Synergien kennen.
Durch den puzzle-artigen Gameplay Loop kommt zusätzlich eine Zielgruppe hinzu, die Freude daran hat, komplexe Spielsituationen zu analysieren, Zusammenhänge zu erkennen und konkrete Lösungswege zu entwickeln.

Eine weitere angrenzende Zielgruppe bilden Spieler etablierter digitaler und analoger Strategie- und Kartenspiele wie \emph{Magic: The Gathering}, \emph{Magic Arena} oder \emph{Hearthstone}.
Die Größenordnung dieses Marktes zeigt sich unter anderem an mehr als 50 Millionen \emph{Magic}-Spielern seit Bestehen der Marke und über zehn Millionen registrierten Spielern von \emph{Magic Arena}.
Auch \emph{Hearthstone} hatte bereits 2018 mehr als 100 Millionen Spieler erreicht.
Diese Spiele unterscheiden sich strukturell deutlich von \coc, sprechen jedoch ebenfalls Spieler an, die Freude am Entdecken und Verstehen komplexer Interaktionen, Kombinationen und Synergien sowie am Aufbau eigener strategischer Spielweisen haben.
Entsprechend gehören auch Spieler moderner strategischer Brett- und Kartenspiele zur erweiterten potenziellen Zielgruppe.

Damit bewegt sich \coc\ in einem bereits etablierten und nachweislich wirtschaftlich attraktiven Markt, ohne ausschließlich von der Zielgruppe eines einzelnen Vergleichstitels abhängig zu sein.
Gleichzeitig zeigt die Bandbreite erfolgreicher Spiele innerhalb dieses Umfelds, dass sich neue Spielkonzepte durch eine klare mechanische Eigenständigkeit innerhalb des Genres positionieren können.

\subsubsection{Unique Selling Points}

\coc\ unterscheidet sich von bisher im Genre etablierten Spielen insbesondere durch folgende Eigenschaften.

\paragraph{Neu entwickelter, eigenständiger Gameplay Loop}
Cycle of Cinder basiert auf einem von uns neu entwickelten, Encounter-basierten Spielprinzip, das sich deutlich von klassischen Kartenkampf- und Deckbuilding-Systemen unterscheidet.

\paragraph{Puzzle-Gameplay mit hoher Player Agency}
Der Spieler löst strategische Problemsituationen mit direkt verfügbaren Fähigkeiten. Der Fokus liegt damit stärker auf Planung, Verständnis und der gezielten Anwendung der eigenen Werkzeuge als auf der Optimierung zufällig gezogener Handkarten.

\paragraph{Ein Core Gameplay Loop, der auch ohne Roguelite-Struktur trägt}
Der zentrale Gameplay Loop funktioniert bereits als eigenständiges Puzzle-Spiel. Im Challenge-Modus wird er in handgefertigten Aufgaben gezielt auf einzelne Mechaniken und Problemstellungen fokussiert. Der Roguelite-Modus erweitert denselben Kern um Build-Entwicklung, Progression und prozedurale Variation. Dadurch entstehen zwei klar voneinander unterscheidbare Spielerlebnisse auf Basis desselben grundlegenden Systems.

\paragraph{Loadout-Building statt klassischem Deckbuilding}
Während eines Runs stellt der Spieler ein Loadout aus Zaubern und Upgrades zusammen, die in den Encountern direkt verfügbar sind. Dadurch bleiben Build-Entscheidungen langfristig relevant und der Zufallseinfluss innerhalb der einzelnen Spielsituationen wird reduziert.

\paragraph{Qualitativ tiefgreifende Upgrade-Systeme}
Upgrades erhöhen nicht lediglich Werte, sondern verändern die Funktionsweise von Fähigkeiten, eröffnen neue Interaktionen und ermöglichen unterschiedliche Spezialisierungen, Synergien und Archetypen.

\paragraph{Roguelite mit prozedural generierten Puzzle-Encountern}
Das Herzstück des Spiels verbindet den puzzle-artigen Gameplay Loop mit prozedural erzeugten Herausforderungen, Run-basierter Build-Entwicklung und persistenter Metaprogression.

\paragraph{Handgefertigter Challenge-Modus}
Zusätzlich zum Roguelite-Modus bietet das Spiel gezielt entwickelte Puzzle-Challenges, in denen jeder Encounter eine qualitativ neue spielerische Fragestellung in den Mittelpunkt stellt.

\paragraph{Community Editor und nutzergenerierte Inhalte}
Spieler können eigene Encounters erstellen und über Steam mit der Community teilen. Dadurch kann der Puzzle-Charakter des Spiels langfristig durch neue, von Spielern entwickelte Herausforderungen erweitert werden.

\paragraph{Fantasy-Welt und fortlaufende Geschichte}
Gameplay, Progression und Encounters sind in eine eigenständige Fantasy-Welt mit Lore, Worldbuilding und einer fortlaufenden Handlung eingebettet.

\subsubsection{Aktueller Entwicklungsstand}
\coc\ ist konzeptionell bereits weit fortgeschritten.
Die grundlegenden Spielsysteme, die Roguelite-Struktur, zentrale Progressionsmechaniken sowie die Spielwelt, ihre Geschichte und die grobe narrative Handlung sind in ihren wesentlichen Grundzügen ausgearbeitet.
Der Schwerpunkt der bisherigen praktischen Entwicklung lag zunächst bewusst auf dem neu entwickelten Gameplay Loop mit puzzle-artigen Levels.
Hierfür wurde ein spielbarer Proof of Concept entwickelt, mit dem sich bereits unabhängig von den späteren Roguelite-Systemen überprüfen lässt, ob das grundlegende Gameplay funktioniert, verständlich ist und Spielern dauerhaft Freude bereitet.
Dieser Schritt war aus unserer Sicht notwendig, bevor weitere Systeme aufgesetzt werden: Die Roguelite-Struktur soll nicht einen noch unbewiesenen Gameplay-Kern interessant machen müssen, sondern auf einem bereits funktionierenden und eigenständig tragfähigen Fundament aufbauen.
Der bestehende Proof of Concept umfasst neben der funktionalen Umsetzung der zentralen Mechaniken bereits eine visuelle Gestaltung, die in wesentlichen Teilen nahe an der für das fertige Spiel angestrebten Präsentation liegt. Damit können neben der grundsätzlichen Funktionsweise des Gameplay Loops bereits heute auch Spielgefühl, Lesbarkeit und visuelle Wirkung unter realistischen Bedingungen getestet werden.

\todo{
	Link zur Demo !!!
}

\subsubsection{Nächster Entwicklungsschritt}
Der nächste wesentliche Entwicklungsschritt besteht darin, das bislang vor allem auf einzelne Level fokussierte Gameplay in einen zunächst vereinfachten Roguelite-Modus zu überführen.
Dadurch sollen erstmals mehrere Level zu einem vollständigen Run verbunden und die zentralen Systeme des späteren Roguelite-Erlebnisses gemeinsam getestet werden.
Dabei müssen insbesondere grundlegende Fragen des Game Balancing und der zeitlichen Struktur beantwortet werden:
Wie viele Level sollte ein Run umfassen? Wie lange sollte ein durchschnittlicher Run dauern? In welchen Abständen benötigt der Spieler Belohnungen? Wie häufig sollten neue Zauber angeboten werden und wie häufig Upgrades bestehender Fähigkeiten? Wie schnell soll sich ein Build innerhalb eines Runs entwickeln und ab welchem Zeitpunkt sollen deutliche Synergien und Archetypen entstehen?
Diese Fragen lassen sich nur eingeschränkt isoliert beantworten. Sie müssen im Zusammenspiel innerhalb eines tatsächlich spielbaren Roguelite-Modus getestet und iterativ angepasst werden.
Auf dieser Basis können anschließend weitere Systeme wie die persistente Metaprogression und die vollständige narrative Einbettung aufgebaut werden.

\subsubsection{Ziel des geförderten Prototyps}
Ziel der Prototypenentwicklung ist eine authentische, spielbare Version von \coc, die den Spielern bereits einen repräsentativen Eindruck davon vermittelt, wie sich das spätere vollständige Spiel anfühlen wird.
Der Prototyp soll daher nicht lediglich einzelne technische Systeme demonstrieren, sondern den neu entwickelten Gameplay-Loop erstmals innerhalb eines zusammenhängenden Roguelite-Erlebnisses spielbar machen.
Dafür soll insbesondere ein Ausschnitt des Roguelite-Modus entstehen, in dem mehrere Level zu einem Run verbunden werden und der Spieler seine Strategie durch neue Zauber und Upgrades kontinuierlich weiterentwickelt.
Gleichzeitig soll die narrative und thematische Einbettung bereits so weit umgesetzt werden, dass Spielwelt, Charaktere und Handlung erkennbar werden und Spieler einschätzen können, welche Rolle diese Elemente im späteren Gesamtprodukt einnehmen.
Der Prototyp ist damit zugleich ein wichtiger Entwicklungsschritt und ein Instrument zur externen Validierung.
Er soll vor dem geplanten Early Access Release unter anderem im Rahmen einer Kickstarter-Kampagne und eines Steam Next Fests öffentlich spielbar gemacht werden.
Potenzielle Spieler sollen das Spiel nicht nur anhand von Beschreibungen, Screenshots oder Videos beurteilen können, sondern selbst erfahren, wie sich der neue Gameplay Loop, die Roguelite-Struktur, das Loadout-Building und die narrative Einbettung zu einem vollständigen Spielerlebnis verbinden.
Das Ziel ist damit ausdrücklich kein isolierter Technologie-Demonstrator, sondern ein repräsentativer Prototyp des späteren Spielerlebnisses, auf dessen Grundlage sowohl das weitere Game Design als auch Balancing, Inhalte und Präsentation bis zum Early Access gezielt weiterentwickelt werden können.
